\section{DrainSinkhorn}
\label{sec:method}

\paragraph{从固定宽度批次到动态压缩。}
标准批量执行并行应用 $W$ 个独立 Sinkhorn 映射，并一直保持宽度 $W$，直到需要最多迭代的问题完成。DrainSinkhorn 以这一固定宽度布局为起点。对问题 $t$，令 $n_t$ 为它首次在已配置验证时点通过双边际验证器时的批量迭代编号。随后，DrainSinkhorn 写出其结果，并动态压缩所有逐问题状态。因此，批量更新 $k$ 按存活宽度
\[
A_k=\left|\{t:n_t\geq k\}\right|
\]
执行。检查间隔是相邻两次验证器检查之间的批量迭代次数。逻辑 masking 在宽度 $W$ 上记录完成情况；压缩后的执行路径则按存活宽度 $A_k$ 运行后续 kernel。固定宽度/动态压缩对照测量释放与动态缩宽的净效应。LM/DC 固定释放决策以隔离动态缩宽。实验分别报告宽度与分组干预。

\subsection{标准批量 Sinkhorn 执行}

令 $F_t$ 表示问题 $t$ 的一次完整 Sinkhorn 循环，其中包含该问题的边缘、kernel 算子和数值表示。标准 batching 增加一个前导问题批次维，但不引入跨问题 reduction。在理想算术下，其代数性质为
\begin{equation}
\widetilde F\!\left(\operatorname{pack}(z_1,\ldots,z_W)\right)
=\operatorname{pack}\!\left(F_1(z_1),\ldots,F_W(z_W)\right).
\label{eq:pack-commutes}
\end{equation}
因此，批量 kernel 保留 EOT 目标和逐实例 Sinkhorn 映射，只改变其执行布局。停止检查后，DrainSinkhorn 过滤问题批次维，并对存活问题应用式 \eqref{eq:pack-commutes}。

对于一个候选，Sinkhorn 循环执行
\begin{equation}
u^{k+1}=a\oslash(Kv^k),\qquad
v^{k+1}=b\oslash(K^\top u^{k+1}),
\label{eq:sinkhorn-cycle}
\end{equation}
其中 $\oslash$ 表示逐元素除法。第二个更新之后，理想算术立即给出
\begin{equation}
v^{k+1}\odot(K^\top u^{k+1})=b.
\label{eq:updated-marginal}
\end{equation}
列边缘由构造保证为最新，而行边缘 $u^{k+1}\odot(Kv^{k+1})$ 仍可能超出容差。DrainSinkhorn 在一次批量操作中计算所有 active 问题的行残差，并且只把预检为正的问题送入双边际验证器。

\begin{figure*}[t]
\centering
\setlength{\fboxsep}{6pt}
\fbox{\parbox{.13\textwidth}{\centering
\textbf{批量更新}\\[-1pt]
$u\leftarrow a/(Kv)$\\
$v\leftarrow b/(K^\top u)$}}
\hspace{2pt}$\rightarrow$\hspace{2pt}
\fbox{\parbox{.13\textwidth}{\centering
\textbf{当前边缘初步检查}\\[-1pt]
批量行残差\\
更新后列恒等式}}
\hspace{2pt}$\rightarrow$\hspace{2pt}
\fbox{\parbox{.13\textwidth}{\centering
\textbf{双边际验证器}\\[-1pt]
当前行与列\\
停止容差}}
\hspace{2pt}$\rightarrow$\hspace{2pt}
\fbox{\parbox{.13\textwidth}{\centering
\textbf{释放并压缩}\\[-1pt]
输出通过的问题\\
压缩全部 OT 状态}}
\hspace{2pt}$\rightarrow$\hspace{2pt}
\fbox{\parbox{.11\textwidth}{\centering
\textbf{下一 kernel 变窄}\\[-1pt]
$W_k\to W_{k+1}$}}
\caption{DrainSinkhorn 的执行循环。使用当前边缘进行初步检查，可以避免许多完整 plan 计算。通过双边际验证器的问题会离开所有逐问题张量；下一次批量更新只处理存活问题。}
\label{fig:active-loop}
\end{figure*}

\subsection{筛查、验证与压缩}

令 $\widehat r^{\mathrm{row}}_{t,k}$ 表示批量预检值。当 $\widehat r^{\mathrm{row}}_{t,k}\leq\tau_{\mathrm{screen}}$ 时，该问题的预检为正；否则它继续执行，而不实例化成本更高的 plan-level 检查。对于预检为正的问题，实现重新计算两侧边缘，并应用配置的停止规则
\begin{equation}
r_{t,k}=\max\!\left\{
\|\pi_{t,k}\ones-a_t\|_1,
\|\pi_{t,k}^\top\ones-b_t\|_1
\right\}\leq\tau_{\mathrm{stop}}.
\label{eq:release-contract}
\end{equation}
被拒绝的问题保持 active，并在后续验证时点再次检查。false-positive 预检会不必要地调用双边际验证器；false-negative 预检会延迟移除。\Cref{sec:experiments} 通过重放筛选结果并逐一使用验证器检查来测量这些效应。

对于每个通过的问题，DrainSinkhorn 写出结果，并压缩所有逐问题对象：source 与 target 描述符、边缘、dual potential、centered shift、log weight 和 scratch buffer。下一次批量更新只看到存活问题。逻辑 mask 对照使用相同的更新、初步检查、双边际验证器、容差和观测完成时间。它冻结通过问题的势函数，但仍在批次张量中保留该问题，因此后续 kernel 仍按原始宽度执行。逻辑 masking 与动态压缩的比较由此隔离缩小批宽所移除的工作。

式 \eqref{eq:pack-commutes} 与式 \eqref{eq:updated-marginal} 定义理想算术下的变换。固定工作 parity、实际残差、阈值附近压力测试和应用端点共同刻画其实现行为。

\paragraph{保持停止规则的释放。}
问题只有通过所选后端使用的同一个双边际验证器后才会被移除。使用当前边缘进行的初步检查只决定是否调用该验证器，并不授权移除。因此，DrainSinkhorn 保留后端已配置的 release authority，只改变未完成问题的调度与存储方式。正式计时的 packed Triton 路径在该后端支持的精度下应用这一规则。独立 FP64 replay 作为有限精度决策行为的验证参照。

\subsection{执行模式与 fallback}

实现提供三种批量执行模式。固定宽度 batching 更新所有问题，直到共同的最终边界。逻辑 masking 记录逐问题完成情况并冻结已完成状态，但不缩窄 kernel。动态压缩移除通过的问题。宽度决定批量 kernel 是否高效，而批量窗口内各问题所需的迭代次数决定动态压缩可以移除多少宽度。正式实验在计时前固定执行模式、宽度、检查间隔和 fallback。正式实验采用 fail-closed 收敛处理：达到迭代上限或出现非有限状态时抛出错误。Packed Triton 实现仅在 dynamic-compaction 路径的存活宽度低于预先配置的 \texttt{min\_packed\_width} 时，才可把窄尾交给匹配的单问题求解器。C66 与 Packer19 正式配置将该阈值设为一。

当一个完整问题能够放入一个加速器时，不同窗口会送往独立的完整设备 worker，从而让每个设备拥有独立更新路径，并避免逐更新 collective。按行划分矩阵仍是单个超大问题的容量路径；主要 scale-out 路径把完整问题分配给独立设备。
